\cabstract{
在学习王老师的大数据原理与技术课程\cite{lecture02}后，在此报告中完成相应的课后作业。

本次作业为“作业-4：电影评论情感分类”。要求基于影评文本实现情感分类，并比较不同方法在准确率与计算效率上的差异。具体要求如下：

\begin{itemize}
    \item 使用 IMDB 影评数据集（或其他可替代数据集）作为实验数据。
    \item 使用 TF-IDF 进行文本特征提取。
    \item 实现至少三类方法中的若干：逻辑回归、支持向量机（SVM）、循环神经网络（RNN）等，并比较它们的分类准确率与计算效率差异。
\end{itemize}

作业提交要求：
\begin{itemize}
    \item 提交内容：代码与简要报告（不少于2页）。
    \item 命名格式：学号+姓名+作业4.zip。
    \item 提交至邮箱：BDSysu2026Spring@163.com。
    \item 截止时间：布置作业之后的下一个周日 24:00。
    \item 严禁抄袭：任何一次抄袭将按零分处理。
\end{itemize}
本实验基于 IMDB 影评数据完成了从数据清洗、TF-IDF 特征构造到多种模型实现与评估的完整流程。主要结论：
在当前的 TF-IDF 表示与实现配置下，线性 SVM是最稳健的选择——它们在准确率和可解释性上均优于本次手写的单步 RNN 与默认配置的 Pegasos SVM（详见第 1 章结果表），并且很好地平衡了计算消耗和预测准确率。
SVM（逐样本 SGD）在修正实现细节后能取得接近 sklearn 基线的性能；而 Pegasos 在本实现及超参数下表现落后，强烈提示需统一正则化定义并在验证集上搜索步长与正则化系数。
序列模型（基于 PyTorch 的 GRU）在序列化文本并使用嵌入后表现优于手写单步 RNN，但仍未超越线性基线，表明在小网格与有限训练轮数下，特征质量（TF-IDF）对最终效果的影响更大。
}
% 中文关键词(每个关键词之间用“，”分开,最后一个关键词不打标点符号。)
\ckeywords{电影评论情感分类，IMDB，TF-IDF，SVM与RNN，情感分析}
% 英文关键词(每个关键词之间用,分开, 最后一个关键词不打标点符号。)
\ekeywords{}


